
\documentclass[UTF8]{ctexart}
%================================ PREAMBLE ==================================

%--------- Packages -----------
\usepackage{fullpage}
\usepackage{amssymb}
\usepackage{amsmath}
\usepackage{amsthm}
\usepackage{latexsym}
\usepackage{graphicx}
\usepackage{wrapfig}
\usepackage{color}
%\usepackage{algorithm,algorithmic}
\usepackage{url}
\usepackage[round]{natbib}
\usepackage[scheme=plain]{ctex}%大写变为小写
\usepackage{lipsum}
\usepackage[colorlinks=true,citecolor=blue]{hyperref}
 
%----------Spacing-------------
%\setlength{\oddsidemargin}{0.25 in}
%\setlength{\evensidemargin}{-0.25 in}
\setlength{\topmargin}{-0.6 in}
%\setlength{\textwidth}{6.5 in}
%\setlength{\textheight}{9.4 in}
\setlength{\headsep}{0.75 in}
\setlength{\parindent}{0 in}
\setlength{\parskip}{0.1 in}

%----------Header--------------
\title{上海市二手房房价影响因素分析研究报告大纲}


\begin{document}
\maketitle%输出标题
\newpage%另起一页
\tableofcontents%命令输出论文目录
\newpage

\section{研究背景} 

\subsection{上海市房地产市场发展现状与行业特征} 

房地产行业是我国国民经济的支柱产业之一，而上海作为我国长三角核心超大城市、全国
房地产市场的标杆城市，其房地产市场运行规律、价格波动特征对全国楼市具有极强的示
范与引导作用。近年来，随着我国房地产市场从高速增量开发转向存量提质发展，上海新
房市场供给严格受限、二手房市场已成为城市住房流通、交易定价的核心主场，二手房房
价走势、空间分布规律，也成为反映上海房地产市场健康程度、城市资源配置效率的核心
指标。

从市场结构来看，上海二手房市场呈现高度微观分化、空间非均衡分布的典型特征。城市
内环、中环、外环间房价梯度差异显著，优质学区、轨道交通、商业医疗配套集中的中心
城区小区房价显著高于远郊区域；同时同一行政区内，相邻小区之间也存在明显的房价联
动、互相影响的特征，传统独立同分布的经典计量模型，已经无法解释上海二手房房价的
微观空间聚集、空间溢出规律。从行业发展特征来看，当前上海楼市已告别普涨周期，政
策端坚持 “房住不炒” 定位、落实稳地价稳房价稳预期目标，市场定价逻辑从单纯板块溢
价，逐步转向小区自身品质、周边公共服务配套、区位通达性等微观因素的综合定价，地
铁通勤、优质教育、生态环境、物业服务等微观配套因素，对二手房房价的影响持续深化。

与此同时，上海城市内部公共服务资源空间分布不均、轨道交通网络持续加密，进一步加
剧了二手房房价的空间关联性：优质配套资源不仅会抬升本小区房价，还会通过空间溢出
效应带动周边相邻小区房价同步上涨，形成高房价聚集片区；而配套薄弱区域则持续形成
低房价聚集片区。这种微观空间单元（小区层面）的房价空间联动性，是上海房地产市场
现阶段最核心的行业特征，也为本研究从微观空间计量视角，开展二手房房价影响因素实
证分析提供了现实基础。

\subsection{二手房房价影响因素的研究意义} 

理论层面，现有二手房房价影响因素研究多基于传统截面回归模型，默认不同小区房价相
互独立，忽略了城市内部微观空间单元的空间相关性与空间溢出效应，无法精准识别配套
因素对房价的直接影响与跨小区溢出效应。本文以上海微观小区二手房为研究对象，引入
空间计量模型、构建地理空间权重矩阵，通过全局 Moran's I 指数检验房价空间聚集特
征，运用空间杜宾模型（SDM）量化各影响因素的直接效应与空间溢出间接效应，弥补传
统计量方法的理论缺陷，完善超大城市二手房微观定价、空间经济学相关研究体系。

现实实践层面，本次研究具备多重应用价值。对购房者而言，可清晰识别地铁、学区、房
龄、物业等核心因素对房价的影响幅度与空间辐射范围，为理性置业选址、房价估值判断
提供科学参考；对房产中介、房企机构而言，能够掌握二手房微观定价逻辑与相邻板块溢
出规律，优化房源估值定价、市场研判策略；对城市规划、住建管理部门而言，可明确公
共服务配套的空间溢价效应，为上海学区均衡化、轨道交通规划、城市公共资源优化配置
、房地产市场平稳调控提供实证数据支撑，助力上海房地产市场健康可持续发展。

\section{数据来源与数据预处理}

\subsection{数据来源与整体说明} 

\subsubsection{核心数据集：上海二手房基础数据集} 

\subsubsection{空间数据：上海市小区经纬度数据} 

\subsection{数据字段与数据概览}

\subsubsection{基础字段说明} 

\subsubsection{经纬度数据与 POI 数据的融合数据} 

\section{上海市二手房房价探索性数据分析}

\subsection{单变量分布特征分析} 

\subsubsection{因变量分析：二手房房价分布特征} 

\subsubsection{连续型自变量分析：面积、楼层、建成年代等分布特征}

\subsubsection{分类型自变量分析：户型、装修情况、房屋朝向等分布特征}

\subsection{相关性分析} 

\subsubsection{房价与建筑属性、小区属性等的相关性分析} 

\subsubsection{变量相关性热力图与多重共线性判断}

\subsection{基于经纬度与 POI 的房价空间分布分析}

\subsubsection{上海市二手房房价整体空间分布格局}

\subsubsection{各行政区 / 板块房价的空间差异与聚类特征}

\subsubsection{核心 POI 对房价的空间影响分析}

\section{实证分析}

\subsection{关键影响变量识别-LASSO 回归模型构建与最优正则化参数选择} 

\subsection{用GLM构建基准模型} 

\subsection{分析基准模型的残差，进行空间效应检验}

\subsection{相同的解释变量构建空间模型}

\subsection{两个模型的解读和拟合效果对比}

\section{结论与总结}

\end{document} 
